\documentclass[11pt,a4paper]{article}
\usepackage{fontspec,amsmath,amssymb,amsthm,microtype}
\setmainfont{lmroman10-regular.otf}[BoldFont=lmroman10-bold.otf,ItalicFont=lmroman10-italic.otf,BoldItalicFont=lmroman10-bolditalic.otf]
\setmonofont{lmmono10-regular.otf}[Scale=MatchLowercase]
\usepackage[margin=25mm]{geometry}
\usepackage[unicode,hidelinks]{hyperref}
\setlength{\emergencystretch}{3em}
\setlength{\parskip}{0.4em}
\setlength{\parindent}{0pt}
\allowdisplaybreaks[2]
\raggedbottom
\title{\texorpdfstring{Quadratic Hilbert identities and integrated convex energy}{Quadratic Hilbert identities and integrated convex energy}}
\author{LangevinLean mathematical exposition}
\date{}
\begin{document}\maketitle
\noindent Source module: \texttt{Proximal\allowbreak{}Langevin/\allowbreak{}Splitting\allowbreak{}L2\allowbreak{}Energy.\allowbreak{}lean}. The accompanying HTML pairs this mathematical exposition with the exact Lean statements and records the scope of the formalization.


The pointwise proximal estimate is lifted to actual \(L^2\) random variables. All scalar energy expectations are proved integrable before distributing integrals or comparing them.


\section{\texorpdfstring{Weighted Young inequality preserves a linear squared contraction}{Weighted Young inequality preserves a linear squared contraction}}


For \(t>0\), expansion and \((t\|v\|-\|r\|)^2\ge0\) prove

\[\|v+r\|^2\le(1+t)\|v\|^2+(1+t^{-1})\|r\|^2.\]

If \(0<t\le1\) and \(\|v\|\le(1-t)\|e\|\), use \((1+t)(1-t)^2\le1-t\). The separate exact identity \(\|v+r+b\|^2=\|v+r\|^2+2\langle v+r+b,b\rangle-\|b\|^2\) keeps an integrated nonsmooth force inside a useful cross term.




Expanding \(\|v+r\|^2\) and using \(2\langle v,r\rangle\le t\|v\|^2+t^{-1}\|r\|^2\) gives the weighted inequality. If \(\|v\|\le(1-t)\|e\|\), \(0<t\le1\), then \((1+t)(1-t)^2\le1-t\). This is the scalar inequality proved in \texttt{hf}; it keeps a linear contraction in the squared error instead of doubling the leading term. The final identity rearranges \(\|v+r+b\|^2\) with its cross term paired against the final error, the form needed for convex-energy cancellation.




\textbf{Correspondence with the displayed declarations.}


\begin{itemize}

\item \texttt{norm\_\allowbreak{}add\_\allowbreak{}sq\_\allowbreak{}le\_\allowbreak{}weighted}: A weighted quadratic triangle bound with an explicit positive weight.
\item \texttt{norm\_\allowbreak{}add\_\allowbreak{}sq\_\allowbreak{}le\_\allowbreak{}of\_\allowbreak{}contraction}: The weighted estimate retains a linear contraction in the squared error.
\item \texttt{norm\_\allowbreak{}sum\_\allowbreak{}rearranged}: Moving an integrated nonsmooth drift into a cross term avoids a Lipschitz constant for that drift in the consistency estimate.
\end{itemize}


{\raggedright\noindent Lean declarations: \texttt{norm\_\allowbreak{}add\_\allowbreak{}sq\_\allowbreak{}le\_\allowbreak{}weighted}, \texttt{norm\_\allowbreak{}add\_\allowbreak{}sq\_\allowbreak{}le\_\allowbreak{}of\_\allowbreak{}contraction}, \texttt{norm\_\allowbreak{}sum\_\allowbreak{}rearranged}.\par}

\section{\texorpdfstring{Justify potential and quadratic expectations}{Justify potential and quadratic expectations}}


For a globally \(K\)-Lipschitz scalar \(f\), center at \(f(0)\); its growth is at most \(K\|X\|\). Hence \(f(X)\in L^2\subset L^1\) on a probability space. The actual Hilbert-space \(L^2\) inner product gives \(\|X\|_2^2=\int\|X(\omega)\|^2dP\), without any arbitrary representative-dependent value.




For a value-Lipschitz observable \(f\), \(|f(X)|\le K\|X\|+|f(0)|\); hence an \(L^2\) state on a probability space gives \(f(X)\in L^2\subset L^1\). The first proof implements this by centering \(f\) at zero. The second identity uses the definition of the \(L^2\) inner product and \(\langle X,X\rangle=\|X\|^2\): \(\|X\|_{L^2}^2=\int\|X(\omega)\|^2dP\). These facts ensure that the energy expressions below are ordinary finite Bochner integrals.




\textbf{Correspondence with the displayed declarations.}


\begin{itemize}

\item \texttt{integrable\_\allowbreak{}lipschitz\_\allowbreak{}comp\allowbreak{}Lp}: Every globally Lipschitz scalar observable of an \texttt{L2} random variable is integrable.
\item \texttt{lp\_\allowbreak{}norm\_\allowbreak{}sq\_\allowbreak{}eq\_\allowbreak{}integral}: Exact second-moment identity for actual \texttt{L2} random variables.
\end{itemize}


{\raggedright\noindent Lean declarations: \texttt{integrable\_\allowbreak{}lipschitz\_\allowbreak{}comp\allowbreak{}Lp}, \texttt{lp\_\allowbreak{}norm\_\allowbreak{}sq\_\allowbreak{}eq\_\allowbreak{}integral}.\par}

\section{\texorpdfstring{Integrate the actual deterministic energy premise}{Integrate the actual deterministic energy premise}}


Suppose \(\|T(x)-y\|^2\le\|x-y\|^2-2h(f(x)-f(y))+C\), with \(T,f\) Lipschitz. Apply it pointwise to \(S,Y\in L^2\). The composition representative agrees almost everywhere with \(T(S)\), and every term is integrable. Monotonicity and linearity of the integral yield

\[\|T(S)-Y\|_2^2\le\|S-Y\|_2^2-2h(\mathbb E f(S)-\mathbb E f(Y))+C.\]


Apply the deterministic premise at \((S(\omega),Y(\omega))\). The representative facts for \texttt{lp\allowbreak{}Drift} and subtraction identify its left side almost everywhere with \(\|T(S)-Y\|^2\). The local facts \texttt{hn}, \texttt{hf\allowbreak{}S}, \texttt{hf\allowbreak{}Y} prove integrability of both quadratic norms and both potential values. Only then does \texttt{integral\_\allowbreak{}mono\_\allowbreak{}ae} integrate the inequality and distribute its finite terms: \[\|T(S)-Y\|_2^2\le\|S-Y\|_2^2-2h(\mathbb Ef(S)-\mathbb Ef(Y))+C.\] The constant remains \(C\) because \(P\) has mass one.




\textbf{Correspondence with the displayed declarations.}


\begin{itemize}

\item \texttt{lp\_\allowbreak{}energy\_\allowbreak{}of\_\allowbreak{}pointwise}: Lift a proved deterministic energy inequality to actual \texttt{L2} random variables. This is an integration lemma; the particle instantiation below proves its pointwise premise from the genuine proximal minimization problem.
\end{itemize}


{\raggedright\noindent Lean declarations: \texttt{lp\_\allowbreak{}energy\_\allowbreak{}of\_\allowbreak{}pointwise}.\par}

\section{\texorpdfstring{Integrate convex supports and split endpoint differences}{Integrate convex supports and split endpoint differences}}


Convex differentiability gives \(\langle S-Y,\nabla e(Y)\rangle\le e(S)-e(Y)\) pointwise. Integrate to obtain the same inequality for the \(L^2\) inner product. For a third variable \(Z\), write \(S-Z=(S-Y)+(Y-Z)\). This splits the cross term into the convex energy difference plus \(\langle Y-Z,\nabla e(Y)\rangle_{L^2}\), whose future-noise part can later be canceled.




The pointwise supporting-plane inequality is \(e(S)-e(Y)\ge\langle\nabla e(Y),S-Y\rangle\). Potential integrability follows from value Lipschitzness; inner-product integrability follows from both vectors belonging to \(L^2\), with score Lipschitzness ensuring \(\nabla e(Y)\in L^2\). Integration gives the first declaration. For an endpoint \(Z\), write \(S-Z=(S-Y)+(Y-Z)\) and use linearity of the inner product. This produces the energy difference plus the reference-increment cross term, precisely the two contributions used in \texttt{splitting\_\allowbreak{}cross\_\allowbreak{}energy\_\allowbreak{}le}.




\textbf{Correspondence with the displayed declarations.}


\begin{itemize}

\item \texttt{lp\_\allowbreak{}convex\_\allowbreak{}gradient\_\allowbreak{}support}: The actual convex gradient supports every square-integrable comparison variable after integrating the pointwise supporting-plane inequality.
\item \texttt{lp\_\allowbreak{}convex\_\allowbreak{}gradient\_\allowbreak{}cross}: Endpoint comparison split into the convex energy difference and a reference-increment cross term, all in the genuine \texttt{L2} inner product.
\end{itemize}


{\raggedright\noindent Lean declarations: \texttt{lp\_\allowbreak{}convex\_\allowbreak{}gradient\_\allowbreak{}support}, \texttt{lp\_\allowbreak{}convex\_\allowbreak{}gradient\_\allowbreak{}cross}.\par}

\end{document}
